\subsubsection{Entendendo endereços}

A memória RAM\footnote{Random Access Memory}, local onde as variáveis são 
guardadas e manipuladas, é essencialmente uma sequência linear de bytes, 
organizada como uma fila, com a capacidade de armazenar informações. Essa 
estrutura pode ser visualizada como uma série de ``caixas'' sequenciais, cada 
uma capaz de reter um valor específico, conforme demonstrado na figura a seguir:

\begin{figure}[htbp]
\centering
\includegraphics[scale=0.2]{memory-address-example}
\label{fig:memoryAdressExample}
\end{figure}

O conceito de endereço refere-se à posição específica de uma variável \(x\) 
dentro desta sequência linear. Esses endereços funcionam como números de 
identificação únicos para a localização de variáveis dentro do espaço de 
memória do programa. Na figura apresentada, os endereços são indicados pelos 
números na parte inferior de cada ``caixa''.

Cada variável vai ocupar uma quantidade diferente de bytes, a depender de seu 
tipo. Por exemplo, variáveis do tipo \texttt{char} ocupam 1 byte, variáveis do 
tipo \texttt{int} ocupam 4 bytes e uma variável do tipo \texttt{double}, 8 
bytes. O tamanho relativo de cada tipo pode ser encontrado na tabela do subcapítulo 
\ref{tipos-de-dados}.

Para acessar o endereço de memória de uma variável, utilizamos o operador \(\&\), conforme exemplificado no código a seguir:

\begin{excode}
    #include <stdio.h>

    int main(void)
    {
        int x = 0;
        printf("Endereço de x: %p\n", &x);

        return 0;
    }
\end{excode}

O código inicializa uma variável \(x\) com o valor 0 e, em seguida, emprega a 
função \texttt{printf}, utilizando o especificador de formato \texttt{\%p} 
para indicar a impressão de um endereço. Dentro da função \texttt{printf}, o 
argumento \texttt{\&x} é utilizado para obter e imprimir o endereço da variável 
\(x\).

Ao executar o código, espera-se um resultado semelhante ao seguinte:

\begin{textbox}
    Endereço de x: 0x7ffd3a450b3c
\end{textbox}

\subsubsection{Bases numéricas}

O valor impresso representa o endereço da variável \(x\) no espaço de memória 
do programa. Nota-se que o endereço é exibido em formato hexadecimal, que 
difere dos números decimais convencionais. A escolha do formato hexadecimal 
(base 16) sobre o decimal é consequência da expansão da capacidade de memória 
ao longo do tempo, tornando os endereços de memória progressivamente maiores e 
mais complexos para representação em base decimal. A base hexadecimal, ao 
permitir uma representação mais compacta de números grandes, tornou-se a 
convenção adotada.

Além do sistema hexadecimal, existem outros sistemas numéricos como o octal 
(base 8) e o binário (base 2) que também podem vir a ser úteis de se conhecer. 
A seguir, apresentamos o número 42.000 representado nessas diferentes bases:

\begin{textbox}
    Hexadecimal: 0xA410
    Octal: 0122020
    Binário: 1010010000010000
\end{textbox}

Observa-se que com o aumento da base numérica, os números representados 
tornam-se mais compactos, enquanto em bases menores, os números expandem. Além 
disso, é importante destacar as convenções para declarar números em diferentes 
bases em C:

\begin{itemize}
        \renewcommand{\labelitemi}{\tiny$\blacksquare$}
    \item Números hexadecimais devem iniciar com \texttt{0x}
    \item Números Octais devem iniciar com \texttt{0}
    \item Números binários devem começar com \texttt{0b}
\end{itemize}

Abaixo, um exemplo prático de como declarar números em diferentes bases na 
linguagem C e imprimir seus valores usando formatações distintas:

\begin{excode}
    #include <stdio.h>

    int main(void)
    {
        int hexa = 0xa410;
        int octal = 0122020; // Não confundir com 122020
        int bin = 0b1010010000010000;

        // %x -> Imprime número em formato hexadecimal
        // %o -> Imprime número em formato octal
        printf("Valor de hexa: %x (%d)\n", hexa, hexa);
        printf("Valor de octal: %o (%d)\n", octal, octal);
        printf("Valor de bin: %d\n", bin);

        return 0;
    }
\end{excode}

Infelizmente, a linguagem C padrão não suporta a formatação direta de números 
binários com \texttt{\%b}. Esta funcionalidade pode requerer implementações 
específicas ou bibliotecas adicionais. Ademais, executando o código acima você 
deve ter o seguinte resultado:

\begin{textbox}
    Valor de hexa: 0xa410 (42000)
    Valor de octal: 122020 (42000)
    Valor de bin: 42000
\end{textbox}
